\chapter{Experimental Setup}\doublespacing
\label{Chapter4}
\lhead{Chapter IV.}

This chapter specifies the empirical environment in which the
framework of Chapter~\ref{Chapter3} is evaluated. It describes
the four datasets used (Section~\ref{sec:datasets}), the
preprocessing pipeline that brings them to a common form
(Section~\ref{sec:preprocessing}), the evaluation metrics
(Section~\ref{sec:metrics}), the implementation details
(Section~\ref{sec:implementation}) and the cross-validation
protocol (Section~\ref{sec:cv}).

\section{Datasets}
\label{sec:datasets}

We use four credit datasets from the UCI Machine Learning
Repository \citep{ucimlrepo}. They were chosen to span a range of
sample sizes, feature counts and class balances; together they
stress the eight classifiers in different ways.
Table~\ref{tab:datasets} summarises their key properties.

\begin{table}[H]
\centering
\caption{Summary of the four UCI credit datasets used in the
benchmark. The class balance is reported as \emph{good payer /
default}.}
\singlespacing
\label{tab:datasets}
\small
\renewcommand{\arraystretch}{1.2}
\begin{tabular}{l c c l l}
\toprule
\textbf{Dataset} & \textbf{Records} & \textbf{Features} & \textbf{Class balance} & \textbf{UCI ID}\\
\midrule
Australian Credit Approval & 690 & 14 & 55.5 / 44.5 (balanced) & 143\\
Japanese Credit Screening & 690 & 15 & 55.5 / 44.5 (balanced) & 27\\
German Credit Data & 1{,}000 & 20 & 70.0 / 30.0 (imbalanced) & 144\\
Taiwan Default of Credit Card Clients & 30{,}000 & 23 & 78.0 / 22.0 (imbalanced) & 350\\
\bottomrule
\end{tabular}
\end{table}
\textbf{Australian Credit Approval.} 690 records, 14 attributes
(6 numeric, 8 categorical), nearly balanced classes. All attribute
names and values are anonymised in the original release to protect
applicant identity.

\textbf{Japanese Credit Screening.} 690 records, 15 attributes
(6 numeric, 9 categorical) and a small number of missing values that
are imputed by mode (categorical) or median (numeric). Like
Australian, it is nearly class-balanced.

\textbf{German Credit Data.} 1{,}000 records, 20 attributes (7
numeric, 13 categorical), 70/30 class imbalance. The accompanying
cost-matrix in the original release weights a missed default five
times more heavily than a missed good applicant; we do not adopt
this cost matrix in our metrics but flag it as a reference for the
relative importance of the two classes.

\textbf{Taiwan Default of Credit Card Clients.} 30{,}000 records,
23 attributes (all numeric: demographics + payment history +
billing amounts), 78/22 class imbalance. It is by far the largest
dataset of the four and is the one on which Bayesian-optimisation
budget and SVM training time are tightest.

\section{Preprocessing}
\label{sec:preprocessing}

Each dataset undergoes the same five-step preprocessing pipeline
before any classifier is trained:

\begin{enumerate}
 \item \textbf{Cleaning.} Drop rows with any remaining missing
 values after standard imputation. Drop the record ID column
 when present.
 \item \textbf{Encoding.} One-hot encode categorical attributes
 using $C-1$ dummy columns for a $C$-level categorical, to avoid
 multicollinearity.
 \item \textbf{Train/test split.} Random 75/25 stratified split
 with random seed 42. Stratification is on the response variable
 so the train and test partitions have the same class proportions.
 \item \textbf{Scaling.} Standardise every feature on the training
 partition only --- subtract the training mean, divide by the
 training standard deviation --- and apply the same transformation
 to the test partition.
 \item \textbf{Imbalance handling.} On the two imbalanced datasets
 (German, Taiwan), apply SMOTE (Equation~\eqref{eq:smote}) to the
 training partition only; otherwise leave the training data as-is.
\end{enumerate}

The choice to apply SMOTE only on the imbalanced datasets is
deliberate. Applying SMOTE on the two balanced datasets
(Australian, Japanese) was found in pilot experiments to lower
$F_1$-macro by 1--2 percentage points, presumably because the
synthetic noise dominates the gain from rebalancing when the data
is already balanced.

\section{Evaluation Metrics}
\label{sec:metrics}

The benchmark reports two \emph{headline} metrics on every
classifier --- $F_1$-macro and AUC --- and a \emph{ten-measure
suite} of diagnostic metrics.

\subsection{Headline Metrics}

\textbf{$F_1$-macro.} Let
$\mathrm{precision}_c$ and $\mathrm{recall}_c$ be the per-class
precision and recall, with $F_1$-score per class
$F_{1,c}=2\,\mathrm{precision}_c \mathrm{recall}_c
/(\mathrm{precision}_c+\mathrm{recall}_c)$. Then
\begin{equation}
F_{1\text{-macro}}
=
\tfrac{1}{2}\bigl(F_{1,0} + F_{1,1}\bigr),
\tag{4.1}
\label{eq:f1macro}
\end{equation}
the unweighted mean of the two per-class $F_1$ scores. Because the
two classes are weighted equally, $F_1$-macro penalises models that
predict the majority class disproportionately, even when accuracy
is high.

\textbf{AUC.} The area under the receiver-operating-characteristic
curve is the probability that the classifier ranks a random positive
above a random negative. It is threshold-independent, so it
isolates the classifier's \emph{ranking} ability from the choice of
operating point.

\subsection{Ten-measure Suite}
\label{sec:ten-measure}

In addition to the headline metrics we report a ten-measure suite,
a diagnostic panel constructed so that no single number can hide a
weakness.  All ten are computed from the entries of the binary
confusion matrix: $TP$ (defaulters correctly identified), $TN$ (good
payers correctly identified), $FP$ (good payers wrongly flagged) and
$FN$ (defaulters wrongly cleared).  Let $N=TP+TN+FP+FN$.  We also
write $\text{sens}=\frac{TP}{TP+FN}$ (sensitivity) and
$\text{spec}=\frac{TN}{TN+FP}$ (specificity).

\textbf{1.\ Accuracy} --- fraction of correct predictions:
\begin{equation}
\text{Accuracy}=\frac{TP+TN}{N}.
\tag{4.2}
\label{eq:acc-tm}
\end{equation}

\textbf{2.\ AUC} --- area under the ROC curve; the probability that
the model ranks a random positive above a random negative.
Threshold-independent.  Defined in Section~\ref{sec:metrics}.

\textbf{3.\ Type-I error rate} --- fraction of negatives wrongly
classified as positive (false-alarm rate):
\begin{equation}
\text{Type-I}=\frac{FP}{FP+TN}=1-\text{spec}.
\tag{4.3}
\label{eq:type1}
\end{equation}

\textbf{4.\ Type-II error rate} --- fraction of positives wrongly
classified as negative (missed-positive rate):
\begin{equation}
\text{Type-II}=\frac{FN}{FN+TP}=1-\text{sens}.
\tag{4.4}
\label{eq:type2}
\end{equation}

\textbf{5.\ EMCC (Expected Misclassification Cost).}  In credit risk
the two errors carry very different consequences: a false negative
(approving a defaulter) loses the loan principal, while a false
positive (rejecting a creditworthy applicant) loses only future
interest revenue.  The standard credit-scoring convention assigns a
cost ratio of $5{:}1$ to FN versus FP.  The expected
misclassification cost is then
\begin{equation}
\text{EMCC}
=
\frac{5\,FN + FP}{N},
\tag{4.5}
\label{eq:emcc}
\end{equation}
the per-test-point expected cost in those units.  Lower is better.

\textbf{6.\ G-Mean.}  The geometric mean of sensitivity and
specificity, balance-sensitive on imbalanced data:
\begin{equation}
\text{G-Mean}
=
\sqrt{\,\text{sens}\cdot\text{spec}\,}
=
\sqrt{\,\frac{TP}{TP+FN}\cdot\frac{TN}{TN+FP}\,}.
\tag{4.6}
\label{eq:gmean}
\end{equation}
G-Mean is zero if either class is completely missed and equals 1 for
a perfect classifier; because it multiplies the two rates rather
than averaging them, it penalises a classifier that excels on one
class but neglects the other.

\textbf{7.\ Discriminant Power (DP).}  Measures how well the
classifier separates the two classes on the score scale:
\begin{equation}
\text{DP}
=
\frac{\sqrt{3}}{\pi}\left[
\log\!\frac{\text{sens}}{1-\text{sens}}
+\log\!\frac{\text{spec}}{1-\text{spec}}
\right].
\tag{4.7}
\label{eq:dp}
\end{equation}
$\text{DP}>3$ indicates strong discriminating power;
$1<\text{DP}<3$ moderate; $\text{DP}<1$ poor.

\textbf{8.\ F-Score.}  The per-class $F_1$ on the default class:
\begin{equation}
\text{F-Score}
=
\frac{2\,\text{precision}\cdot\text{recall}}{\text{precision}+\text{recall}}
=
\frac{2TP}{2TP+FP+FN},
\tag{4.8}
\label{eq:fscore}
\end{equation}
with $\text{precision}=\frac{TP}{TP+FP}$ and
$\text{recall}=\text{sens}$.  The headline $F_1$-macro of
Section~\ref{sec:metrics} averages F-Score over both classes.

\textbf{9.\ Cohen's $\kappa$.}  Accuracy corrected for chance
agreement:
\begin{equation}
\kappa
=
\frac{\text{Accuracy}-p_e}{1-p_e},
\quad
p_e=\frac{(TP+FP)(TP+FN)+(FN+TN)(FP+TN)}{N^{2}},
\tag{4.9}
\label{eq:kappa}
\end{equation}
where $p_e$ is the probability of a chance agreement.
$\kappa = 1$ for perfect agreement, $\kappa = 0$ for chance-level
agreement.  Cohen's $\kappa$ is also the metric that the GHOST
procedure uses internally to choose the decision threshold
(Equation~\eqref{eq:ghost}).

\textbf{10.\ Youden's J.}  A single-number summary that rewards
high sensitivity and high specificity equally:
\begin{equation}
J=\text{sens}+\text{spec}-1.
\tag{4.10}
\label{eq:youden}
\end{equation}
$J$ ranges from $-1$ to $1$; $J=0$ means random-guessing performance
on the chosen class, $J=1$ a perfect classifier on both classes.

The ten-measure suite is computed on every (dataset,
feature-selection, classifier) row; the per-row results are stored
in the \texttt{<dataset>\_v3\_ten\_measure\_results.csv} files used
as inputs to Chapter~\ref{Chapter5}.

\section{Implementation}
\label{sec:implementation}

All experiments are implemented in Python~3.10. The key libraries
are: \texttt{scikit-learn} 1.3 for the base learners and the
cross-validation utilities; \texttt{xgboost} 2.0 for XGBoost;
\texttt{imbalanced-learn} 0.11 for SMOTE; \texttt{scikit-optimize} 0.9
for Bayesian optimisation; \texttt{ghostml} 0.4 for GHOST threshold
tuning; \texttt{scipy.stats} for the Friedman and Wilcoxon tests;
and \texttt{ucimlrepo} 0.0.7 for the data fetch.

The GA-MaxAccLogit meta-learner is implemented as a custom
scikit-learn-compatible estimator in pure Python so it can be
slotted into the same cross-validation and pipeline machinery as
the other meta candidates. The implementation follows
Algorithm~\ref{alg:gamal} exactly and is reproduced in
Appendix~\ref{AppendixA}.

A fixed random seed (42) is used throughout for reproducibility,
including for the random train/test split, the SMOTE oversampler,
the cross-validation fold assignment, the GHOST threshold scan
and the GA-MaxAccLogit population initialisation. The only
remaining source of non-determinism is the floating-point order of
operations in some multi-threaded library calls; this affects the
last decimal of some metrics on the larger Taiwan dataset but does
not change any model ranking.

\section{Cross-validation Protocol}
\label{sec:cv}

The benchmark distinguishes between three nested cross-validations,
each serving a different role:

\textbf{CV-1 (Bayesian-optimisation tuning).} Inside the training
partition, a $k$-fold cross-validation is used to score each
hyperparameter setting tried by the Bayesian optimiser. $k=5$ on
Australian/Japanese/German and $k=3$ on Taiwan (to fit the
computational budget).

\textbf{CV-2 (Stacker out-of-fold prediction).} Also inside the
training partition: a stratified 5-fold cross-validation generates
out-of-fold predictions for each base learner. These predictions
form the meta-feature matrix $\mathbf{Z}$ used to train the meta
candidates. This is the leakage-safe step of
Section~\ref{sec:stacking}.

\textbf{Out-sample test set.} The 25\% test partition that was set
aside at the beginning is used \emph{once}, at the very end, to
score every classifier. All Bayesian-optimisation tuning,
GHOST threshold selection and best-meta selection happen on the
training data; the test partition is not touched until evaluation.
This protocol guarantees that the reported $F_1$-macro and AUC are
\emph{generalisation} estimates rather than in-sample fits.

The benchmark therefore produces one row of metrics per (dataset,
feature-selection, classifier) combination, for a total of
$4\times 4\times 8 = 128$ rows. These rows form the input to the
statistical analysis of Chapter~\ref{Chapter5}.
